\section{X : Optimising Execution Speed (\& Memory)}
\label{chap:optim}

%%%%%%%%%%%%%%%%%%%%%%%%%%%%%%%%%%%%%%%%%%%%%%%%%%%%%%%%%%%%%%

\begin{frame}[fragile]
\frametitle{Bitwise Operators}
\begin{columns}[T]

\begin{column}{0.40\textwidth}
\lstinputlisting[style=basicc]{../Code/ChapX/twiddle.c}
\end{column}

\pause
\begin{column}{0.50\textwidth}
\begin{itemize}[<+->]
\item For some operations bitwise manipulation makes
sense.
\item Use \verb^>>^, \verb^<<^, \verb^&^ and \verb^|^.
\item However, you might believe that shifting an integer right one bit
might be quicker than dividing by two.
\item Inspect the assembled code, OR switch off optimisation \verb^-O0^ and time a large
loop performing each of the approaches.
\item On a modern processor with a modern compiler, it probably makes
no difference, so use the approach that is most self-documenting.
\end{itemize}
\end{column}

\end{columns}
\end{frame}

%%%%%%%%%%%%%%%%%%%%%%%%%%%%%%%%%%%%%%%%%%%%%%%%%%%%%%%%%%%%%%

\begin{frame}[fragile]
\frametitle{Memory Alignment : Structures}
\begin{columns}[T]

\begin{column}{0.45\textwidth}
Why can you not directly compare two structures?

\lstinputlisting[style=basicc]{../Code/ChapX/structcomp.c}
\end{column}

\pause
\begin{column}{0.45\textwidth}
\begin{itemize}[<+->]
\item The reason is somewhat complex.
\item The elements of a structure are memory aligned.
\item This allows the processor to quickly access elements
without complicated (expensive) bit-shifting.
\item The structure may contain padding (slop) whose contents are undefined.
\item The compiler optimises for execution speed over total memory used.
\end{itemize}
\end{column}

\end{columns}
\end{frame}

%%%%%%%%%%%%%%%%%%%%%%%%%%%%%%%%%%%%%%%%%%%%%%%%%%%%%%%%%%%%%%


\begin{frame}[fragile]
\frametitle{Memory Alignment : Structures}
\begin{columns}[T]

\begin{column}{0.50\textwidth}
\lstinputlisting[style=basicc]{../Code/ChapX/alin.c}
\end{column}

\pause
\begin{column}{0.40\textwidth}
\begin{itemize}[<+->]
\item Running this produces:
\begin{verbatim}
0  4  8 = 12
0  1  5 = 7
\end{verbatim}
\item The compiler has quietly added `slop' between elements in the first structure.
\item Therefore, using \verb^memcmp()^ to compare between two structures of the same type might fail.
\item The second structure has no `slop' between elements.
\item Put the largest elements first in the \verb^struct^, and the smallest last.
\item \url{http://www.catb.org/esr/structure-packing/}
\end{itemize}
\end{column}

\end{columns}
\end{frame}

%%%%%%%%%%%%%%%%%%%%%%%%%%%%%%%%%%%%%%%%%%%%%%%%%%%%%%%%%%%%%%

\begin{frame}[fragile]
\frametitle{Memory Alignment : SIMD}
\begin{columns}[T]


\pause
\begin{column}{0.50\textwidth}
\begin{itemize}[<+->]
\item Most modern CPUs have instruction set extensions which allow the same operation to be performed on multiple data {\bf simultaneously}.
\item This is know as SIMD (Single Instruction Multiple Data).
\item Here we want to sum two arrays element-wise.
\end{itemize}
\end{column}

\pause
\begin{column}{0.40\textwidth}
\lstinputlisting[style=basicc]{../Code/ChapX/avx.c}
\end{column}

\end{columns}
\end{frame}

%%%%%%%%%%%%%%%%%%%%%%%%%%%%%%%%%%%%%%%%%%%%%%%%%%%%%%%%%%%%%%

\begin{frame}[fragile]
\frametitle{Memory Alignment : SIMD}
\begin{columns}[T]

\begin{column}{0.60\textwidth}
\lstinputlisting[style=basicc]{../Code/ChapX/avx2.c}
\end{column}

\pause
\begin{column}{0.35\textwidth}
\begin{itemize}[<+->]
\item Here we use AVX2 (available since 2013)
\item We need a \verb^C11^ compiler (\verb^-std=c11^).
\item \verb^float^s are $32$-bits, and registers are $256$-bits.
\item Notice we're stepping through the additions {\bf eight} at a time.
\item Does this make the code eight times quicker though?
\item An optimizing compiler (\verb^-O3^) may do some of this heavy lifting for you anyway.
\end{itemize}
\end{column}

\end{columns}
\end{frame}

%%%%%%%%%%%%%%%%%%%%%%%%%%%%%%%%%%%%%%%%%%%%%%%%%%%%%%%%%%%%%%

\begin{frame}[fragile]
\frametitle{Multi-Processors}
\begin{columns}[T]

\pause
\begin{column}{0.45\textwidth}
\begin{itemize}[<+->]
\item OpenMP allows some simple additions to your C code
which enable Multi-Processor execution.
\item You takes chunks of your code (typically a big loop)
and share them out over several processors/threads.
\item Use flag \verb^-fopenmp^ to compile.
\end{itemize}
\end{column}

\pause
\begin{column}{0.45\textwidth}
\lstinputlisting[style=basicc]{../Code/ChapX/thread.c}
\end{column}


\end{columns}
\end{frame}

%%%%%%%%%%%%%%%%%%%%%%%%%%%%%%%%%%%%%%%%%%%%%%%%%%%%%%%%%%%%%%

\begin{frame}[fragile]
\frametitle{OpenMP}
\begin{columns}[T]

\pause
\begin{column}{0.45\textwidth}
\begin{itemize}[<+->]
\item Here we have a big loop that goes through all possibilities
and counts the numbers of primes.
\end{itemize}
\end{column}

\pause
\begin{column}{0.45\textwidth}
\lstinputlisting[style=basicc]{../Code/ChapX/primes_no.c}
\end{column}


\end{columns}
\end{frame}

%%%%%%%%%%%%%%%%%%%%%%%%%%%%%%%%%%%%%%%%%%%%%%%%%%%%%%%%%%%%%

\begin{frame}[fragile]
\frametitle{OpenMP}
\begin{columns}[T]

\pause
\begin{column}{0.50\textwidth}
\lstinputlisting[style=basicc,linerange={1-25}]{../Code/ChapX/primes.c}
\end{column}

\pause
\begin{column}{0.40\textwidth}
\begin{itemize}[<+->]
\item Now each loop is farmed-out to different processors.
\item {\bf BUT}: The variable \verb^cnt^ is over-written by all the
threads independently.
\item We have to make a copy of \verb^cnt^ that is local to each thread
and sum them together at the end (a reduction).
\item There are diferent types of scheduling (which iterations of the loop
is assigned to which processor): \verb^static, dynamic, guided, auto,^ and \verb^runtime^.
\end{itemize}
\end{column}

\end{columns}
\end{frame}
